% !TeX root = ../tfg.tex
% !TeX encoding = utf8
%
% ====================================================================
%  APÉNDICE A — PREGUNTAS REALIZADAS A LA EMPRESA Y RESPUESTAS
%
%  RECORDATORIO AL EQUIPO:
%  - Este apéndice es OBLIGATORIO según las instrucciones de la asignatura.
%  - Debéis incluir las preguntas que realizasteis en la entrevista/visita
%    a la empresa y las respuestas que os dieron.
%  - Organizad las preguntas por temática (Tema 2 y Tema 3) o por el
%    orden en que se realizaron.
%  - Podéis incluir el nombre o cargo del entrevistado (si da permiso)
%    y la fecha y medio de la entrevista (presencial, videollamada...).
%  - Si hay información que el entrevistado pidió que no se publicara,
%    indicadlo (p.ej. "Dato confidencial — no se puede reproducir").
%  - Si grabasteis la entrevista (con permiso), podéis transcribir los
%    fragmentos más relevantes.
% ====================================================================

\chapter{Entrevista a CaixaBank: preguntas y respuestas}

% ============================================================
\section{Datos de la entrevista}
% ============================================================

\begin{table}[H]
  \centering
  \begin{tabularx}{\textwidth}{lX}
    \toprule
    \textbf{Parámetro}       & \textbf{Detalle} \\
    \midrule
    Empresa                  & CaixaBank \\
    Nombre del entrevistado  & [Anónimo a petición] \\
    Cargo                    & [Director de Recursos Humanos regional] \\
    Departamento             & [Recursos Humanos] \\
    Fecha                    & [12/04/2026] \\
    Duración                 & [15 min] \\
    Formato                  & [Telefónica] \\
    \bottomrule
  \end{tabularx}
\end{table}

% ============================================================
\section{Preguntas sobre el Tema 4 — Proceso de reclutamiento}
% ============================================================

\subsubsection*{¿Cuáles son los costes directos e indirectos asociados habitualmente al proceso de reclutamiento en la empresa? ¿Existe un presupuesto limitado para esto?}

\textbf{Respuesta:}\newline
``(El entrevistado sólo supo que acudían a una ETT para parte de este proceso).''

\subsubsection*{¿Cuál es el tiempo medio estimado (Time to Fill) desde que se detecta la necesidad de cubrir un puesto hasta que el candidato se incorpora?}

\textbf{Respuesta:}\newline
``(No supo responder el entrevistado).''

\subsubsection*{¿Qué enfoque de reclutamiento predomina actualmente en la organización: interno, externo o mixto?}

\textbf{Respuesta:}\newline
``Mixto. Normalmente el proceso de una vacante comienza en un portal interno en el cuál prevalece por un periodo de tiempo estipulado para cada tipo de vacante, de forma que predomina el ``moving'' para personal interno, ofreciendo así la posibilidad a los compañeros de adquirir nuevas responsabilidades y cambiar su ámbito de trabajo a personas que realmente cumplan los requisitos para la misma. También se valoran diferentes aspectos del proyecto en el que la persona esté colaborando actualmente para conocer si es posible la reducción de ``recursos'' asignados, se valorarán las cualidades del empleado para llevar a cabo un estudio de la aportación de conocimientos del ámbito, soft skills y hard skills, para elaborar una puntuación sobre la eficiencia de la persona en ambos proyectos para así poder optar por un bien estar para el empleado, además de sacar partido a su ``expertise''.''

\subsubsection*{A nivel interno: ¿Qué canales o métodos se utilizan para comunicar las vacantes a los empleados actuales (intranet, tablón de anuncios, programas de referidos)?}

\textbf{Respuesta:}\newline
``A nivel interno la forma de comunicar las vacantes disponibles a partir de la propia intranet ya mencionada.''

\subsubsection*{A nivel externo: ¿Cuáles son las fuentes de reclutamiento externo más efectivas para la empresa (portales de empleo, agencias de colocación, universidades, ferias de empleo)?}

\textbf{Respuesta:}\newline
``Dependiendo de la campaña utilizan diferentes medios de publicación, adaptando los canales de comunicación al tipo de perfil que buscan y al objetivo estratégico. Por ejemplo, si se trata de atraer talento joven para programas específicos como el CaixaBank Talent Program, es habitual reforzar la presencia en universidades, ferias de empleo y campañas digitales segmentadas en redes sociales profesionales.

En cambio, para perfiles más senior o especializados, se priorizan portales como LinkedIn o colaboraciones con consultoras de selección y headhunters. Asimismo, plataformas como InfoJobs o Glassdoor permiten dar visibilidad masiva a las vacantes.

De esta forma, la estrategia de reclutamiento externo no es única, sino que se ajusta en función del tipo de campaña, el público objetivo y la urgencia o especialización del perfil requerido, combinando distintos canales para maximizar el alcance y la calidad de las candidaturas.''

\subsubsection*{¿Se usan nuevas tendencias de reclutamiento, como por ejemplo 2.0 (Linkedin, Twitter, Facebook para atraer candidatos), 3.0 (employer branding) o video cv?}

\textbf{Respuesta:}\newline
``Actualmente, la mayoría de empresas utilizan una combinación de herramientas automatizadas e inteligencia artificial para filtrar grandes volúmenes de candidaturas de forma rápida y eficiente. 

La técnica más utilizada es el uso de sistemas ATS (Applicant Tracking System), que actúan como un primer filtro automatizado. Estos sistemas analizan los currículums, extraen la información y clasifican a los candidatos en función de su adecuación al puesto. También se usa machine learning para comparar a los candidatos con las ofertas a las que aplican y predecir su encaje, además de otras aplicaciones de la IA como análisis predictivos o automatización de distintos procesos como programar entrevistas.

Video cv en principio no, pero sí aceptan cartas de recomendación.''

\subsubsection*{¿Qué nuevas estrategias, tecnologías o tendencias tiene en mente aplicar la empresa en un futuro próximo para optimizar la captación de talento?}

\textbf{Respuesta:}\newline
``(No se le preguntó al entrevistado.)''

% ============================================================
\section{Preguntas sobre el Tema 5 — Proceso de selección e integración}
% ============================================================

\subsubsection*{Preselección: ¿Qué criterios (filtros killer, palabras clave, experiencia) se utilizan para realizar la criba curricular inicial?}

\textbf{Respuesta:}\newline
``Para la criba inicial se hace uso de sistemas ATS (Applicant Tracking System), que tienen unos filtros automatizados para cribar currículums. Estos filtros se basan en palabras clave, y se tiene en cuenta tanto filtros ``killer'' (requisitos excluyentes como nivel de idioma o formación específica dependiendo el puesto), experiencia profesional y coherencia del CV (progresión, estabilidad, adecuación al puesto), siendo este el criterio más importante, y la ubicación geográfica y disponibilidad.''

\subsubsection*{Pruebas: ¿Cómo se determina qué tipo de pruebas de selección se aplicarán a cada perfil profesional?}

\textbf{Respuesta:}\newline
``Dependiendo del tipo de perfil profesional o del nivel de responsabilidad del puesto. A perfiles técnicos se les suele plantear un caso práctico o real, mientras que a perfiles junior más pruebas de razonamiento para ver su potencial.''

\subsubsection*{Decisión y contratación: ¿Quiénes intervienen en la toma de decisión final (RRHH, responsable de departamento) y cómo se comunica la oferta al candidato elegido?}

\textbf{Respuesta:}\newline
``Intervienen el departamento de RRHH junto con los directores de área. Comunica la decisión recursos humanos para puestos bajos, o el director y oficial de RRHH para puestos más altos.''

\subsubsection*{¿En qué casos específicos se aplican pruebas escritas u orales (psicotécnicos, conocimientos, idiomas) y cómo se evalúan?}

\textbf{Respuesta:}\newline
``(El entrevistado sabe de la realización de pruebas escritas psicotécnicas en el pasado, pero desconoce qué ocurre en la realidad).''

\subsubsection*{¿Qué estructura siguen las entrevistas individuales (estructuradas, semi-estructuradas, por competencias) y qué rol juegan las entrevistas colectivas o dinámicas de grupo?}

\textbf{Respuesta:}\newline
``Son entrevistas estructuradas y se realizan dinámicas de grupo en ocasiones.''

\subsubsection*{¿Se utilizan técnicas de simulación (Role-play, método del caso) o Centros de Evaluación (Assessment Centers) para puestos directivos o clave?}

\textbf{Respuesta:}\newline
``Sí se usan.''

\subsubsection*{¿Se han implementado las vídeo entrevistas (ya sean en directo vía videollamada o en diferido/grabadas) en las primeras fases del proceso?}

\textbf{Respuesta:}\newline
``Se realizan video entrevistas, pero no suelen realizarse en las primeras fases del proceso, sino más adelante.''

\subsubsection*{Formal: ¿Qué acciones programadas (cursos de formación, reuniones de bienvenida, presentación de normativas) componen la socialización del nuevo empleado?}

\textbf{Respuesta:}\newline
``Sí hay, tanto para nuevos como antiguos. Tanto presenciales como online, cursos, reuniones, programas de bienvenida, una formación inicial obligatoria, etc.''

\subsubsection*{Informal: ¿Cómo se fomenta la integración social y cultural del empleado en el día a día (desayunos, asignación de un ``compañero'' o mentor no oficial, eventos de teambuilding)?}

\textbf{Respuesta:}\newline
``Se realizan tanto eventos sociales, como desayunos, se utilizan las mentorías para nuevos empleados, y eventos de teambuilding.''

\subsubsection*{¿Dispone la organización de un Manual de Acogida estructurado y actualizado?}

\textbf{Respuesta:}\newline
``Sí existe, mencionando principalmente los valores de la entidad, así como las funciones concretas del puesto. Sí está actualizado.''

\subsubsection*{¿Qué métricas o feedback se recogen de los nuevos empleados para evaluar la eficacia del proceso de selección y acogida?}

\textbf{Respuesta:}\newline
``Se realizan evaluaciones periódicas (6 meses, 1 año) mencionando qué cualidades destacan para el individuo en el puesto, como aquellas que son mejorables.''

\subsubsection*{¿Qué innovaciones tecnológicas (ej. Inteligencia Artificial para la preselección, gamificación en pruebas) o cambios metodológicos planea incorporar la empresa próximamente en esta área?}

\textbf{Respuesta:}\newline
``IA sí, gamificación también se está implementando.''

